\section*{Capítulo \ref{ch:b3:cancer-biology} --- Biologia do câncer}

\begin{solution}{exo:b3:cancer-biology:1}
Um \omterm{def:b3:cancer-biology:genes}{oncogene} é uma versão com ganho de função de um gene que promove a
proliferação ou a sobrevivência; um alelo mutante basta (dominante):
\emph{RAS} (mutação pontual), \emph{MYC} (translocação, amplificação),
\emph{HER2}, \emph{BCR--ABL}. Um supressor de \omterm{def:b3:cancer-biology:hallmarks}{tumor} contém a
proliferação ou impõe parada, morte ou reparo; os dois alelos precisam
ser perdidos (recessivo no nível celular): \emph{RB}, \emph{TP53},
\emph{APC}, \emph{BRCA1}.
\end{solution}

\begin{solution}{exo:b3:cancer-biology:2}
Proliferação sustentada: \emph{RAS}, \emph{MYC}. Insensibilidade aos
supressores de crescimento: \emph{RB}, \emph{CDKN2A} (p16). Resistência
à morte: \emph{BCL2}, \emph{TP53}. Imortalidade replicativa: o promotor
da \omterm{def:b3:cancer-biology:telomerase}{telomerase}. Angiogênese: \emph{VEGF} induzido pelo HIF (perda de
\emph{VHL}). \omterm{def:b3:cancer-biology:metastasis}{Invasão} e \omterm{def:b3:cancer-biology:hallmarks}{metástase}: perda da E-caderina, os fatores da
\omterm{def:b3:cancer-biology:metastasis}{transição epitélio--mesênquima}. Instabilidade do genoma: genes de
\omterm{prop:b3:dna-repair:fidelity}{reparo de mau pareamento}, \emph{\omterm{def:b3:dna-repair:dsb}{BRCA}}. Metabolismo desregulado:
\emph{MYC}, \omterm{prop:b3:cancer-biology:pathways}{via PI3K}. \omterm{prop:b3:cancer-biology:immune}{Evasão imune}: \emph{PD-L1}, perda de MHC.
Inflamação que favorece o \omterm{def:b3:cancer-biology:hallmarks}{tumor}: sinalização por
\emph{NF-$\kappa$B}.
\end{solution}

\begin{solution}{exo:b3:cancer-biology:3}
As crianças com história familiar desenvolviam vários \omterm{def:b3:cancer-biology:hallmarks}{tumores} nos dois
olhos, cedo; os casos esporádicos desenvolviam um único \omterm{def:b3:cancer-biology:hallmarks}{tumor}, mais
tarde. Ele inferiu que dois eventos eram necessários numa célula: os
pacientes familiares haviam herdado um e precisavam só de um segundo
evento somático (frequente o bastante para ocorrer em várias células da
retina), ao passo que os esporádicos precisavam que os dois ocorressem
somaticamente numa mesma célula --- raro, tardio, único. Daí um gene
cujas duas cópias precisam ambas ser perdidas: o primeiro supressor de
\omterm{def:b3:cancer-biology:hallmarks}{tumor}.
\end{solution}

\begin{solution}{exo:b3:cancer-biology:4}
Um mutante de \emph{Salmonella} que exige histidina é plaqueado sem
histidina junto com o composto em teste; só surgem colônias de células
em que uma nova mutação reverteu o defeito, de modo que o número delas
mede a mutagenicidade. Acrescenta-se extrato de fígado porque muitos
carcinógenos são inofensivos até que as enzimas hepáticas os convertam
em metabólitos reativos, como acontece no corpo; as bactérias não têm
essas enzimas.
\end{solution}

\begin{solution}{exo:b3:cancer-biology:5}
A incidência sobe $100$ vezes enquanto a idade dobra: $2^{k-1} = 100$,
$k - 1 =
\log_{2}100 = 6.6$, ou seja, cerca de sete ou oito eventos
limitantes da taxa na leitura ingênua --- menos, se a expansão clonal
entre os golpes for admitida.
\end{solution}

\begin{solution}{exo:b3:cancer-biology:6}
$L = \sqrt{2Dc_{0}/q}$: para $q = 0.01$, $\sqrt{2\times 2\times 10^{-9}
\times 0.05/0.01} = \qty{141}{\micro m}$; para $q = 0.03$,
\qty{82}{\micro m}. Um \omterm{def:b3:cancer-biology:hallmarks}{tumor} de crescimento rápido respira mais
depressa, esgota o oxigênio a uma distância menor, e seu centro morre
num tamanho menor.
\end{solution}

\begin{solution}{exo:b3:cancer-biology:7}
$10^{11}\times 10^{-3n} < 1$ exige $n > 3.7$: quatro ciclos. Depois de
dois ciclos restam $10^{5}$ células --- um décimo de milímetro cúbico,
muito abaixo dos $10^{9}$ que se conseguem detectar ---, de modo que o
\omterm{def:b3:cancer-biology:hallmarks}{tumor} está ``sumido'' por qualquer teste enquanto cem mil células
permanecem; o tratamento precisa continuar pelo número planejado de
ciclos, e não até o \omterm{def:b3:cancer-biology:hallmarks}{tumor} desaparecer.
\end{solution}

\begin{solution}{exo:b3:cancer-biology:8}
As duas lesões ativam a mesma via (receptor $\to$ \omterm{def:b3:cancer-biology:genes}{Ras} $\to$ ERK); uma
vez presente uma delas, a outra não confere vantagem adicional e não é
selecionada. Um inibidor de EGFR bloqueia o receptor; uma célula que
adquira uma mutação em KRAS a jusante restaura o sinal com o receptor
ainda bloqueado, e é selecionada durante o tratamento --- resistência
por desvio.
\end{solution}

\begin{solution}{exo:b3:cancer-biology:9}
A E7 libera o E2F da Rb, levando a célula epitelial infectada pelo
\omterm{def:b3:cell-cycle-apoptosis:checkpoints}{ponto de restrição} sem fatores de crescimento; a E6 destrói a \omterm{def:b3:cancer-biology:genes}{p53}, de
modo que a proliferação inapropriada e o dano que ela causa não
disparam nem parada nem \omterm{def:b3:cell-cycle-apoptosis:apoptosis}{apoptose}. Duas marcas são fornecidas de uma só
vez, mas as restantes --- imortalidade, \omterm{def:b3:cancer-biology:metastasis}{invasão}, instabilidade ---
precisam de mais mutações somáticas, que levam décadas para se acumular
nas células persistentemente infectadas. A vacina suscita anticorpos
neutralizantes contra o capsídeo e previne a infecção; dada antes da
exposição, nunca há uma célula infectada a transformar.
\end{solution}

\begin{solution}{exo:b3:cancer-biology:10}
Depois do $i$-ésimo golpe o clone cresce por $g$, de modo que $g^{i}$
vezes mais células ficam sob risco do golpe seguinte; a probabilidade
da sequência completa na linhagem descendente de uma célula original é
$g\cdot g^{2}\cdots$ --- para um $g$ constante por golpe, o risco
cumulativo se torna $N(ut)^{k}
g^{k-1}$ a menos de fatores dependentes
da ordem, e a incidência mantém a potência $t^{k-1}$ com um prefator
$g^{k-1}$ maior. Se as próprias expansões crescerem com o tempo (clones
em crescimento exponencial), a curva fica ainda mais íngreme, de modo
que uma inclinação de $6$ é produzida por menos de sete condutoras: o
$k$ de Armitage--Doll é um limite superior do número de condutoras
limitantes da taxa, e os \omterm{def:b3:cancer-biology:hallmarks}{tumores} sequenciados mostram de duas a oito.
\end{solution}

\begin{solution}{exo:b3:cancer-biology:11}
(a) O fármaco A sozinho enfrenta $\mu_{A}N = 100$ células resistentes
preexistentes: $P_{0} = e^{-100} \approx 0$; o clone resistente a A
sobrevive e volta a crescer durante as dezoito semanas de A, e, quando
B começa, ele já é de novo uma população grande, de modo que B também
enfrenta $\mu_{B}N' \gg 1$ --- a sequência falha duas vezes. (b)
Juntos: uma célula duplamente resistente surge a $10^{-14}$ por
divisão, $\mu_{A}\mu_{B}N = 10^{-5}$, $P_{0} =
0.99999$. A combinação
desde o início, quando $N$ é menor, é a regra porque é o produto de
duas taxas pequenas que torna a resistência improvável.
\end{solution}

\begin{solution}{exo:b3:cancer-biology:12}
Com $100$ vezes mais células e os mesmos $u$ e $k$, o risco cumulativo
$N(ut)^{k}$ seria $100$ vezes maior --- todo elefante deveria morrer de
\omterm{def:b3:cancer-biology:hallmarks}{câncer} jovem. Não morre, logo $u$ ou $k$ precisam diferir: vinte cópias
de \emph{TP53} tornam mais forte a resposta apoptótica ao dano e
removem células mutantes antes que elas se expandam (baixando o $u$
efetivo daquela etapa); os animais grandes também têm taxas de mutação
menores por célula (melhor reparo, menos divisões por célula por ano,
porque suas células se renovam mais devagar) e podem exigir mais golpes
(uma camada extra de supressores). A supressão do \omterm{def:b3:cancer-biology:hallmarks}{câncer} evoluiu junto
com o tamanho corporal.
\end{solution}

\begin{solution}{pb:b3:cancer-biology:1}
\textbf{1.} $10^{12}/10^{3} = 10^{9}$ células; $\log_{2}10^{9} \approx 30$
duplicações.
\textbf{2.} $30\times 100 = 3000$ dias, cerca de $8$ anos.
\textbf{3.} \qty{1}{kg} $\approx 10^{12}$ células: mais dez
duplicações, $1000$ dias, menos de três anos --- um terço do tempo até
a detecção.
\textbf{4.} O tempo de duplicação no início precisa ter sido de cerca de
$60$ dias, ou o \omterm{def:b3:cancer-biology:hallmarks}{tumor} cresceu mais depressa quando pequeno e desacelerou
à medida que crescia (o padrão usual, à medida que o suprimento e o
espaço falham).
\textbf{5.} Um ciclo de $2$ dias dobraria a população a cada $2$ dias;
um tempo de duplicação de $100$ significa que o crescimento líquido é
$2\,\%$ da taxa de nascimento: quase toda célula que nasce é compensada
por uma que morre, ou poucas estão ciclando --- nascimento e morte
quase em equilíbrio, e o crescimento é a pequena diferença.
\textbf{6.} Os fármacos matam células em ciclo, de modo que um \omterm{def:b3:cancer-biology:hallmarks}{tumor} com
uma fração pequena em ciclo perde uma fração menor por ciclo; mas o
\omterm{def:b3:cancer-biology:hallmarks}{tumor} rápido, tendo perdido mais, também volta a crescer entre os
ciclos. No balanço, são os \omterm{def:b3:cancer-biology:hallmarks}{tumores} rápidos (leucemias, linfomas, \omterm{def:b3:cancer-biology:hallmarks}{câncer}
de testículo) os que a quimioterapia cura.
\textbf{7.} Razão $300$ para uma razão de idades $80/30 = 2.67$:
$k - 1 =
\ln 300/\ln 2.67 = 5.8$, $k \approx 7$.
\textbf{8.} $N(ut)^{k} = 10^{8}\times (7\times 10^{-5})^{7} \approx
10^{-21}$: absurdo contra um risco ao longo da vida perto de $1$ em
$3$. Falta: a expansão clonal que multiplica as células sob risco
depois de cada golpe, as taxas de mutação elevadas pela instabilidade,
muito mais divisões (células-tronco ciclando por décadas) e menos
condutoras verdadeiras.
\textbf{9.} Multiplicar por $g^{k-1} = 100^{6} = 10^{12}$ dá $10^{-9}$
--- ainda ínfimo com sete eventos; com quatro eventos e expansões,
$10^{8}\times (7\times 10^{-5})^{4}\times 10^{6} \approx 2\times
10^{-3}$, a ordem certa. O quadro verdadeiro é de poucas condutoras e
grandes expansões.
\textbf{10.} $(ut)^{k-1}/(ut)^{k} = 1/(ut) = 1/(10^{-6}\times 40) =
2.5\times 10^{4}$: a incidência de uma portadora aos quarenta é dezenas
de milhares de vezes a esporádica --- os \omterm{def:b3:cancer-biology:hallmarks}{cânceres} hereditários atacam
cedo.
\textbf{11.} Um evento dez vezes mais rápido: incidência $\times 10$;
dois: $\times 100$. O risco observado, de quinze a vinte e cinco vezes,
sugere que a fumaça age em mais de uma etapa.
\textbf{12.} Ao parar, o $u$ volta a cair para os eventos que ainda
faltam, e por isso a incidência deixa de subir tão abruptamente; mas os
golpes já acumulados nas células permanecem, de modo que as células do
ex-fumante estão mais perto do \omterm{def:b3:cancer-biology:hallmarks}{câncer} que as de quem nunca fumou, e o
excesso de risco persiste, congelado em vez de apagado.
\textbf{13.} $L = \sqrt{2\times 2\times 10^{-9}\times 0.05/0.03} =
\qty{82}{\micro m}$.
\textbf{14.} Totalmente oxigenado até um diâmetro $2L \approx
\qty{160}{\micro m}$. Um esferoide de \qty{2}{mm} tem uma borda viva de
\qty{82}{\micro m}: fração viva $1 - (918/1000)^{3} = 0.23$.
\textbf{15.} As células a meio caminho estão a \qty{75}{\micro m} de um
vaso, logo dentro de $L$ --- oxigenadas em princípio, mas os vasos
tumorais carregam menos oxigênio e têm fluxo intermitente, de modo que
as células do meio ficam hipóxicas. As células hipóxicas são duas a três
vezes mais resistentes à radiação (é preciso oxigênio para tornar o dano
permanente), e a hipóxia dispara \omterm{def:b3:cell-cycle-apoptosis:apoptosis}{apoptose} dependente de \omterm{def:b3:cancer-biology:genes}{p53}, de modo que
seleciona as células que perderam a \omterm{def:b3:cancer-biology:genes}{p53}.
\textbf{16.} Reduzir a densidade à metade afasta os vasos por
$\sqrt{2}$: \qty{212}{\micro m} de distância, com as células do meio a
\qty{106}{\micro m}, além de $L$: elas morrem ou, sobrevivendo
hipóxicas, tornam-se mais agressivas e angiogênicas.
\textbf{17.} $30/2 = 15$ vezes mais glicose. Os \omterm{def:b3:cancer-biology:hallmarks}{tumores}, por isso,
concentram análogos de glicose, e uma tomografia com
fluordesoxiglicose emissora de pósitrons os acende.
\textbf{18.} \omterm{def:b3:cancer-biology:hallmarks}{Tumores} famintos encolhem até o tamanho limitado pela
difusão e persistem, dormentes e hipóxicos, ou cooptam vasos já
existentes; a hipóxia seleciona clones agressivos. Mas podar a
vasculatura caótica do \omterm{def:b3:cancer-biology:hallmarks}{tumor} baixa a pressão intersticial e uniformiza o
fluxo, de modo que a quimioterapia penetra melhor --- a combinação
funciona onde nenhuma das duas funciona sozinha.
\textbf{19.} $10^{9}\times 10^{-2n} < 1$: $n > 4.5$, cinco ciclos.
\textbf{20.} $21$ dias $= 0.21$ duplicação, recrescimento $\times 1.16$;
fator líquido por ciclo $0.0116$; $\ln 10^{-9}/\ln 0.0116 = 4.65$: ainda
cinco ciclos.
\textbf{21.} $10^{-7}\times 10^{9} = 100$ células resistentes esperadas;
$P_{0} = e^{-100} \approx 0$: a resistência é certa.
\textbf{22.} $\mu_{1}\mu_{2}N = 10^{-14}\times 10^{9} = 10^{-5}$;
$P_{0} = 0.99999$. Com $10^{12}$ células, $10^{-2}$, $P_{0} = 0.99$.
\textbf{23.} $10^{6}$ células: um fármaco $\mu N = 0.1$, $P_{0} = 0.90$;
dois fármacos $10^{-8}$, $P_{0} \approx 1$. Trate quando a população é
pequena, e trate com combinações.
\textbf{24.} Os linfócitos~T atacam muitos neoantígenos ao mesmo tempo,
de modo que nenhuma mutação isolada escapa deles; escapar exige perder a
própria apresentação de antígenos (MHC,
$\beta_{2}$-microglobulina) ou suprimir os linfócitos~T, uma classe de
evento diferente e mais rara. A mesma carga alta de mutações que torna
um \omterm{def:b3:cancer-biology:hallmarks}{tumor} fadado a resistir a um inibidor de cinase lhe dá muitos
neoantígenos: os \omterm{def:b3:cancer-biology:hallmarks}{tumores} deficientes em \omterm{prop:b3:dna-repair:fidelity}{reparo de mau pareamento} e os
induzidos por ultravioleta e por fumaça respondem melhor.
\textbf{25.} Cerca de $8$ anos até a detecção; \omterm{def:b3:cancer-biology:hallmarks}{tumor} vivo até
\qty{82}{\micro m} em torno de um vaso; $P_{0} = 0.99999$ de que uma
combinação de dois fármacos não enfrente célula resistente alguma num
\omterm{def:b3:cancer-biology:hallmarks}{tumor} de \qty{1}{cm^3}.
\end{solution}
